\documentclass[11pt]{amsart}

\usepackage[T1]{fontenc}
\usepackage{lmodern}
\usepackage{microtype}
\usepackage{mathtools,amssymb,amsthm}
\usepackage[hidelinks]{hyperref}

\hypersetup{
  pdftitle={A 13-Vertex Graph With No 4-Cycle Whose Worst Orientation Beats the Matching Bound}
}

\newtheorem{theorem}{Theorem}
\newtheorem{lemma}[theorem]{Lemma}
\newtheorem{proposition}[theorem]{Proposition}
\theoremstyle{definition}
\newtheorem*{problem}{Open problem 37 of Bousquet, Deschamps, Lehtil\"a and Parreau}
\theoremstyle{remark}
\newtheorem{remark}[theorem]{Remark}

\DeclareMathOperator{\mold}{mold}
\newcommand{\gammaLD}{\gamma_{\mathrm{LD}}}
\newcommand{\Nin}{N^-}

\title[A graph with no $4$-cycle beating the matching bound]
{A 13-Vertex Graph With No $4$-Cycle Whose Worst Orientation\\ Beats the Matching Bound}

\date{}

\subjclass[2020]{05C69, 05C20, 05C35}
\keywords{locating-dominating set, oriented graph, worst orientation,
matching number, graph without $C_4$, open problem}

\begin{document}

\begin{abstract}
For a graph $G$ let $\mold(G)$ be the largest value of the locating-domination
number $\gammaLD(D)$ over all orientations $D$ of $G$. Bousquet, Deschamps,
Lehtil\"a and Parreau proved that $\mold(G)\le n-\alpha'(G)$ whenever $G$ has
no $C_4$ subgraph, where $n=|V(G)|$ and $\alpha'$ is the matching number, and
asked whether the inequality is ever strict (Open problem 37 in the e-print of
their paper). It is. We exhibit a graph $G$ on $13$ vertices and $22$ edges
with no $4$-cycle, graph6 string \texttt{L??E@agXAp@wdC}, for which
\[
 \mold(G)=6<7=n-\alpha'(G).
\]
The lower bound $\mold(G)\ge6$ is exhibited by one orientation and is checkable
by hand; the upper bound $\mold(G)\le6$ is a finite statement about all
$2^{22}=4\,194\,304$ orientations. We do not claim that $G$ is smallest or
unique.
\end{abstract}

\maketitle

\section{The problem}

All graphs are finite, simple and undirected; an \emph{orientation} $D$ of $G$
assigns to each edge one of its two directions. For $u\in V(D)$ write
$\Nin_D(u)$ for the set of in-neighbours of $u$. Following Bousquet,
Deschamps, Lehtil\"a and Parreau \cite{BDLP}, a set $S\subseteq V(D)$ is
\emph{locating-dominating} in the oriented graph $D$ when the sets
\[
 I_D(u)=\Nin_D(u)\cap S \qquad (u\in V(D)\setminus S)
\]
are all non-empty and pairwise distinct; $\gammaLD(D)$ denotes the least size
of such a set, which exists because $S=V(D)$ is vacuously
locating-dominating. The parameter studied in \cite[Section~4]{BDLP} is the
\emph{worst} orientation value
\[
 \mold(G)=\max\{\,\gammaLD(D)\;:\;D\text{ an orientation of }G\,\}.
\]
Say that $G$ is \emph{without $C_4$} if $G$ has no $4$-cycle as a subgraph,
induced or not; equivalently, no two vertices of $G$ -- adjacent or not -- have
two common neighbours. This is the hypothesis of \cite[Lemma~34]{BDLP}, which
states that
\begin{equation}\label{eq:L34}
 \mold(G)\;\le\;n-\alpha'(G)
 \qquad\text{for every graph $G$ of order $n$ without $C_4$,}
\end{equation}
$\alpha'(G)$ being the matching number. Immediately after
\cite[Corollary~36]{BDLP} they ask:

\begin{problem}
Does there exist a graph $G$ without $C_4$ as a subgraph with
$\mold(G)<n-\alpha'(G)$?
\end{problem}

\noindent The statement is quoted from the e-print \texttt{arXiv:2112.01910v2}
of \cite{BDLP}, where it is numbered Open problem~37. In
\cite{BDLP} ``without $H$'' and ``$H$-free'' are different, the latter being
reserved for the induced version, so the hypothesis of the problem is that $G$
contains no $4$-cycle at all.

\begin{theorem}\label{thm:main}
Let $G$ be the graph on $V=\{0,1,\dots,12\}$ with graph6 string
\texttt{L??E@agXAp@wdC}, that is, with the $22$ edges
\[
\begin{array}{llllll}
 \{0,6\},&\{0,8\},&\{0,12\},&\{1,6\},&\{1,9\},&\{1,10\},\\
 \{2,7\},&\{2,8\},&\{2,9\},&\{3,7\},&\{3,10\},&\{3,12\},\\
 \{4,8\},&\{4,10\},&\{4,11\},&\{5,9\},&\{5,11\},&\{5,12\},\\
 \{6,11\},&\{7,11\},&\{8,10\},&\{9,12\}. & &
\end{array}
\]
Then $G$ has no $C_4$ subgraph, $n=13$, $\alpha'(G)=6$ and $\mold(G)=6$. In
particular
\[
 \mold(G)=6<7=n-\alpha'(G),
\]
so Open problem 37 of \cite{BDLP} has an affirmative answer.
\end{theorem}

\section{The graph, and its matching number}\label{sec:obj}

The adjacency lists of $G$ are
\[
\begin{array}{llll}
 0:6,8,12 & 1:6,9,10 & 2:7,8,9 & 3:7,10,12\\
 4:8,10,11 & 5:9,11,12 & 6:0,1,11 & 7:2,3,11\\
 8:0,2,4,10 & 9:1,2,5,12 & 10:1,3,4,8 & 11:4,5,6,7\\
 12:0,3,5,9 & & &
\end{array}
\]
so $G$ is connected with eight vertices of degree $3$, five of degree $4$, and
$\sum_v\deg v=44=2\cdot22$. The edge list above is written in lexicographic
order and we keep that order throughout: an \emph{orientation index} is an
integer $o$ with $0\le o<2^{22}$, and bit $i$ of $o$ is $1$ exactly when the
$i$-th listed edge $\{a,b\}$ (with $a<b$) is oriented $a\to b$. Thus bit $0$
governs the edge $\{0,6\}$ and bit $21$ the edge $\{9,12\}$.

\begin{lemma}\label{lem:noC4}
$G$ has no $C_4$ subgraph.
\end{lemma}

\begin{proof}
A $4$-cycle $uxvy$ consists of two distinct vertices $u,v$ together with two
distinct common neighbours $x,y$ of them, and conversely any two distinct
vertices with two distinct common neighbours span a $4$-cycle. The pair $u,v$
need \emph{not} be adjacent, so the condition to check is that every one of the
$\binom{13}{2}=78$ unordered pairs of vertices has at most one common
neighbour, not merely each of the $22$ adjacent ones. Inspection of the
thirteen adjacency lists above shows that it holds: $54$ of the $78$ pairs have
exactly one common neighbour and the other $24$ have none, in agreement with
the count $\sum_v\binom{\deg v}{2}=8\binom32+5\binom42=54$ of paths of length
two.
\end{proof}

\begin{lemma}\label{lem:alpha}
$\alpha'(G)=6$, hence $n-\alpha'(G)=7$; and $A=\{0,1,2,3,4,5\}$ is independent
in $G$.
\end{lemma}

\begin{proof}
The six edges $\{0,6\},\{1,9\},\{2,7\},\{3,10\},\{4,11\},\{5,12\}$ are listed
above and have twelve distinct endpoints, so they form a matching and
$\alpha'(G)\ge6$. Since $n=13$ is odd, $\alpha'(G)\le\lfloor13/2\rfloor=6$.
Hence $\alpha'(G)=6$ and $n-\alpha'(G)=13-6=7$. Finally, none of the fifteen
pairs inside $A$ appears in the edge list, so $A$ is independent.
\end{proof}

\section{$\mold(G)\ge6$, from a single orientation}\label{sec:lower}

\begin{proposition}\label{prop:lower}
$\mold(G)\ge6$.
\end{proposition}

\begin{proof}
Let $A=\{0,1,2,3,4,5\}$, independent by Lemma~\ref{lem:alpha}. Exactly $18$ of
the $22$ edges have an endpoint in $A$ -- all of them but $\{6,11\}$,
$\{7,11\}$, $\{8,10\}$ and $\{9,12\}$ -- and each of those $18$ has its other endpoint in
$B=\{6,\dots,12\}$. Let $D^{*}$ be any orientation of $G$ in which all $18$ of
those edges are directed \emph{away} from $A$; the remaining four edges, which
lie inside $B$, may be oriented arbitrarily, so there are $16$ such
orientations and the argument below applies to each.

In $D^{*}$ every vertex of $A$ is a source, i.e.\ $\Nin_{D^{*}}(u)=\emptyset$
for $u\in A$. Hence if $S$ is locating-dominating in $D^{*}$ and $u\in A$ were
outside $S$, then $I_{D^{*}}(u)=\emptyset$, which is forbidden. So
$A\subseteq S$ and $\gammaLD(D^{*})\ge|A|=6$.

Conversely $S=A$ is locating-dominating in $D^{*}$: the seven traces are
\[
 I(6)=\{0,1\},\quad I(7)=\{2,3\},\quad I(8)=\{0,2,4\},\quad I(9)=\{1,2,5\},
\]
\[
 I(10)=\{1,3,4\},\quad I(11)=\{4,5\},\quad I(12)=\{0,3,5\},
\]
that is, $I(v)=N(v)\cap A$ for each $v\in B$. They are non-empty, and they are
pairwise distinct: three of them have size~$2$ and four have size~$3$, and the
three of size~$2$ are pairwise distinct, as are the four of size~$3$. Hence
$\gammaLD(D^{*})=6$ and $\mold(G)\ge6$.
\end{proof}

\section{$\mold(G)\le6$: the one finite computation}\label{sec:upper}

\begin{lemma}[monotonicity]\label{lem:mono}
Let $D$ be an orientation of a graph and let $S\subseteq S'\subseteq V(D)$. If
$S$ is locating-dominating in $D$, so is $S'$.
\end{lemma}

\begin{proof}
Let $v\in V(D)\setminus S'$. Then $v\notin S$ and
$I'(v):=\Nin_D(v)\cap S'\supseteq\Nin_D(v)\cap S=I(v)$, which is non-empty, so
$I'(v)\neq\emptyset$. Moreover $S\subseteq S'$ gives $I'(v)\cap S=I(v)$ for every
such $v$; hence if $I'(u)=I'(v)$ for two distinct $u,v\notin S'$ then
$I(u)=I'(u)\cap S=I'(v)\cap S=I(v)$, contradicting that $S$ is
locating-dominating. No hypothesis on the graph is used.
\end{proof}

In particular, for $k\le n$ the inequality $\gammaLD(D)\le k$ holds if and only
if some set of size \emph{exactly} $k$ is locating-dominating in $D$, so it is
enough to examine sets of one size.

\begin{proposition}\label{prop:upper}
Every one of the $2^{22}=4\,194\,304$ orientations of $G$ admits a
locating-dominating set of size $6$. Hence $\mold(G)\le6$, and with
Proposition~\ref{prop:lower}, $\mold(G)=6$.
\end{proposition}

This is a finite statement, but not hand-sized: the number of orientations
times the $\binom{13}{6}=1716$ candidate six-sets is about $7\times10^{9}$
pairs. Here is the reduction used.

Fix a $6$-set $S$ and put $T=V\setminus S$, $|T|=7$. Whether $S$ is
locating-dominating in an orientation $D$ depends on $D$ only through the seven
traces $I_D(v)=\Nin_D(v)\cap S$, $v\in T$; and $I_D(v)$ is determined by the
directions of the edges joining $v$ to $S$. Those edge sets, over $v\in T$, are
pairwise disjoint, and edges inside $S$ or inside $T$ are irrelevant. Inside
the $2^{22}$-element space of orientation indices:
\begin{itemize}
\item[(a)] ``$I_D(v)=\emptyset$'' is the set of indices agreeing with one
  prescribed pattern on the $|N(v)\cap S|$ bits of the edges from $v$ to $S$,
  and free on the other $22-|N(v)\cap S|$ bits;
\item[(b)] ``$I_D(u)=I_D(v)$'' for $u\ne v$ in $T$ is the union, over the
  non-empty subsets $c\subseteq N(u)\cap N(v)\cap S$, of the sets of indices
  agreeing with one prescribed pattern on the bits of the edges from $u$ to $S$
  and from $v$ to $S$. By Lemma~\ref{lem:noC4} the set $N(u)\cap N(v)$ has at
  most one element, so there is at most one such $c$ and the union has at most
  one term.
\end{itemize}
Each of the sets in (a) and (b) is a subcube of the orientation cube: the
indices agreeing with a prescribed pattern on a set $F$ of bit positions are
obtained from the single index carrying that pattern by the $22-|F|$ doublings
$X\mapsto X\cup(X+2^{q})$, one per free position $q$, each a shift-and-or on a
$2^{22}$-bit word. So for fixed $S$ the set of orientations in which $S$ is
\emph{not} locating-dominating is a union of at most $7+\binom72=28$ such
subcubes; unioning the complements over the six-sets $S$ and comparing with the
all-ones mask decides Proposition~\ref{prop:upper} exactly.

Carrying this out, $980$ of the $1716$ six-sets already suffice: the union of
their masks is all of $\{0,1,\dots,2^{22}-1\}$. Proposition~\ref{prop:upper}
follows, and hence Theorem~\ref{thm:main}.

\section{Scope}\label{sec:status}

Existence is all that is claimed. We do not claim that $G$ is of least order or
least size, nor that it is unique, and nothing here shows that $13$ is the
least order at which the answer is affirmative. The bound
$\mold(G)\le n-\alpha'(G)$ of \cite[Lemma~34]{BDLP} remains a theorem; we show
only that it is not always tight, and we do not determine for which graphs
without $C_4$ it is tight.

Neighbouring settings: locating-dominating sets in graphs of
girth at least $5$ \cite{BFH}, a subclass of the graphs without $C_4$ that does
not contain $G$, since $G$ has the triangle $\{4,8,10\}$; location-domination
together with matchings in undirected cubic graphs \cite{FH}; the
\emph{minimum} of the identifying-code number over all orientations \cite{CH};
best and worst orientations for the metric dimension \cite{ABC}; and $\gammaLD$
of local tournaments \cite{BBLP}.

Lemmas~\ref{lem:noC4} and \ref{lem:alpha} and Proposition~\ref{prop:lower} need
no machine. A program that carries out the computation of
Proposition~\ref{prop:upper} from the graph6 string alone by the reduction of
Section~\ref{sec:upper}, together with a recorded run of it, accompanies this
note; the program also re-derives the statements proved above by hand and
cross-checks the mask construction against the naive per-orientation predicate.

\begin{thebibliography}{9}

\bibitem{BDLP}
N.~Bousquet, Q.~Deschamps, T.~Lehtil\"a, and A.~Parreau,
\emph{Locating-dominating sets: from graphs to oriented graphs},
Discrete Math. \textbf{346} (2023), 113124,
\href{https://doi.org/10.1016/j.disc.2022.113124}{doi:10.1016/j.disc.2022.113124}.
Statement labels and numbering are quoted from the e-print
\href{https://arxiv.org/abs/2112.01910}{arXiv:2112.01910v2}.

\bibitem{ABC}
J.~Araujo, J.~Bensmail, V.~Campos, F.~Havet, A.~K. Maia, N.~Nisse, and
A.~Silva,
\emph{On finding the best and worst orientations for the metric dimension},
Algorithmica (2023),
\href{https://doi.org/10.1007/s00453-023-01132-0}{doi:10.1007/s00453-023-01132-0}.

\bibitem{BBLP}
T.~Bellitto, C.~Brosse, B.~Lev\^eque, and A.~Parreau,
\emph{Locating-dominating sets in local tournaments},
Discrete Appl. Math. (2023),
\href{https://arxiv.org/abs/2109.03102}{arXiv:2109.03102}.

\bibitem{BFH}
C.~Balbuena, F.~Foucaud, and A.~Hansberg,
\emph{Locating-dominating sets and identifying codes in graphs of girth at
least 5},
Electron. J. Combin. \textbf{22} (2015), no.~2, Paper 2.15.
\href{https://arxiv.org/abs/1407.7263}{arXiv:1407.7263}.

\bibitem{CH}
N.~Cohen and F.~Havet,
\emph{On the minimum size of an identifying code over all orientations of a
graph},
Electron. J. Combin. \textbf{25} (2018), no.~1, Paper 1.49.

\bibitem{FH}
F.~Foucaud and M.~A. Henning,
\emph{Location-domination and matching in cubic graphs},
Discrete Math. \textbf{339} (2016), no.~4, 1221--1231.
\href{https://arxiv.org/abs/1412.2865}{arXiv:1412.2865}.

\end{thebibliography}

\end{document}
